\chapter{T\'opicos de An\'alisis N\'umerico}\label{ch:AnalNum}

\section{Ecuaciones nolineales}

\textcite{heath2002scientificcomputing}

\section{Optimizaci\'on}

\subsection{}

\section{M\'etodos de integraci\'on n\'umerica de EDOs}

\subsection{Paso temporal fijo}

\subsection{Paso temporal variable}

\section{Descomposici\'on de valores singulares (SVD)}

Referencia: \textcite[Cap. 1]{brunton_data-driven_2019}

Dado $X \in \mathbb{C}^{n \times m}$:

\begin{equation}
    X = U \Sigma V^*
\end{equation}

\noindent donde $U  \in \mathbb{C}^{n \times n} $ y $V \in \mathbb{C}^{m \times m}$ son matrices unitarias con columnas ortogonales (es decir, $U U^* = I$ y $V V^* = I$), y $\Sigma \in \mathbb{R}^{n \times m}$ es una matriz con entradas no-negativas y reales en la diagonal, y ceros fuera de la diagonal. Las columnas de $U$ se denominan los \emph{vectores singulares izquierdos} de $X$, y las columnas de $V$ son los \emph{vectores singulares derechos} de $X$.

Cuando $n \geq m$, se puede escribir la operación SVD económica:

\begin{equation}
    X = U \Sigma V^* =
    \begin{bmatrix}
        \hat{U} & \hat{U}^{\perp}    
    \end{bmatrix}
    \begin{bmatrix}
        \hat{\Sigma} \\ 0    
    \end{bmatrix}
    V^*
\end{equation}

donde $\hat{\Sigma} \in \mathbb{R}^{m \times m}$ es la matriz de valores singulares.

En Matlab, se utiliza la función \texttt{svd()}

\begin{lstlisting}[style=Matlab-editor,frame=single,numbers=left]
    X = rand(5,3);
    [U,S,V] = svd(X);
    [Uhat,Shat,V] = svd(X,'econ');
\end{lstlisting}

Una interpretación matemática es que las matrices de correlación $X X^*$ y $X X^*$ cumple las siguientes relaciones:

\begin{align}
    (X X^*) U & = U \begin{bmatrix} \hat{\Sigma}^2 & 0 \\ 0 & 0\end{bmatrix}\\
    (X^* X) V & = V \hat{\Sigma}^2\\
\end{align}

Entonces, los valores singulares de $X$ corresponden a los valores propios (\emph{eigenvalues}) de las matrices de correlación $X X^*$ y $X X^*$. También, las columnas de $U$ son los vectores propios (\emph{eigenvectors}) de la matriz $X X^*$, y las columnas de $V$ son los vectores propios de la matriz $X^* X$.

\section{Fast Fourier Transform}

